\chapter{\IfLanguageName{dutch}{Resultaten}{Results}}%
\label{ch:resultaten}

Dit hoofdstuk bespreekt de resultaten van de experimenten uit
Hoofdstuk~\ref{ch:experimentopzet}. Eerst wordt toegelicht hoe de
uitvoeringstijden werden gemeten. Daarbij wordt een onderscheid gemaakt tussen
scenario's met één virtuele machine en scenario's met meerdere virtuele machines.
Vervolgens worden de metingen voor Linux en Windows weergegeven. Ten slotte
worden de resultaten van RustProv en Vagrant met elkaar vergeleken.

\section{Meetmethode}
\label{sec:meetmethode}

Om de uitvoeringstijden op een consistente manier te registreren, werden twee
PowerShell-scripts gebruikt. Het eerste script is bedoeld voor scenario's met
één virtuele machine. Het tweede script wordt gebruikt voor scenario's met
meerdere virtuele machines. Beide scripts meten de duur van een opgegeven
opdracht en schrijven het resultaat weg naar een CSV-bestand.

Bij één virtuele machine worden de starttijd en de provisioningtijd afzonderlijk
gemeten. De starttijd is de tijd die nodig is om de VM op te starten. De
provisioningtijd is de tijd die nodig is om het provisioningsscript uit te
voeren.

Bij meerdere virtuele machines wordt alleen de totale uitvoeringstijd van de
volledige opstelling gemeten. Deze tijd omvat zowel het opstarten als het
provisionen van alle VM's. De VM's worden daarbij na elkaar verwerkt.

Elke meting wordt drie keer uitgevoerd onder dezelfde omstandigheden. Hierdoor
wordt vermeden dat één uitzonderlijk snelle of trage uitvoering het resultaat
bepaalt. Op basis van de drie afzonderlijke metingen wordt per scenario het
gemiddelde berekend. Dit gemiddelde wordt gebruikt om RustProv en Vagrant met
elkaar te vergelijken.

Voor één virtuele machine wordt de totale uitvoeringstijd berekend als de som van
de gemiddelde starttijd en de gemiddelde provisioningtijd:

\[
\overline{T}_{\mathrm{totaal}} =
\overline{T}_{\mathrm{start}} +
\overline{T}_{\mathrm{provisioning}}
\]

Voor meerdere virtuele machines wordt de totale uitvoeringstijd rechtstreeks
gemeten vanaf het starten van de volledige opstelling tot het moment waarop alle
VM's succesvol zijn geprovisioned.

Dezelfde meetmethode wordt toegepast op RustProv en Vagrant. De gebruikte
virtuele machines, configuraties, provisioningscripts en
hardware-instellingen blijven daarbij gelijk.

\subsection{PowerShell-meetscripts}

Voor de meting van één virtuele machine wordt een afzonderlijk PowerShell-script
gebruikt. Dit script meet de starttijd en de provisioningtijd afzonderlijk. Het
volledige script wordt weergegeven in
Listing~\ref{lst:meet-script-one-vm}.

\begin{lstlisting}[language=bash,
    caption={PowerShell-script voor het meten van één virtuele machine.},
    label={lst:meet-script-one-vm}]
    % Het PowerShell-meetscript voor één VM wordt hier ingevoegd.
\end{lstlisting}

Voor scenario's met meerdere virtuele machines wordt een tweede PowerShell-script
gebruikt. Dit script meet enkel de totale uitvoeringstijd van de volledige
opstelling. Het volledige script wordt weergegeven in
Listing~\ref{lst:meet-script-multiple-vms}.

\begin{lstlisting}[language=bash,
    caption={PowerShell-script voor het meten van meerdere virtuele machines.},
    label={lst:meet-script-multiple-vms}]
    % Het PowerShell-meetscript voor meerdere VM's wordt hier ingevoegd.
\end{lstlisting}

Beide scripts registreren per meting minstens de gebruikte tool, het scenario en
de uitvoeringstijd. De tijd wordt uitgedrukt in seconden. De resultaten worden
opgeslagen in een CSV-bestand, zodat de afzonderlijke metingen kunnen worden
verwerkt en het gemiddelde per scenario kan worden berekend.

\section{Meting van één Linux-VM}

Voor het eerste scenario werd één AlmaLinux-VM gebruikt. De starttijd en de
provisioningtijd werden afzonderlijk gemeten voor zowel RustProv als Vagrant. De
VM werd voor elke nieuwe meting in dezelfde beginsituatie geplaatst.

Bij RustProv werd de starttijd gemeten met het commando \texttt{start}. Nadat de
VM was opgestart, werd de provisioningtijd gemeten met het commando
\texttt{provision}. Bij Vagrant werd dezelfde procedure gevolgd met
respectievelijk \texttt{vagrant up} en \texttt{vagrant provision}.

\begin{table}[htbp]
    \centering
    \caption{Meting van één Linux-VM met RustProv en Vagrant.}
    \label{tab:meting-een-linux}
    \begin{tabular}{llrrrr}
        \toprule
        \textbf{Tool} &
        \textbf{Onderdeel} &
        \textbf{Run 1 (s)} &
        \textbf{Run 2 (s)} &
        \textbf{Run 3 (s)} &
        \textbf{Gemiddelde (s)} \\
        \midrule
        RustProv & Start & -- & -- & -- & -- \\
        RustProv & Provisioning & -- & -- & -- & -- \\
        Vagrant  & Start & -- & -- & -- & -- \\
        Vagrant  & Provisioning & -- & -- & -- & -- \\
        \bottomrule
    \end{tabular}
\end{table}

Tabel~\ref{tab:meting-een-linux} toont de afzonderlijke metingen en de
gemiddelde uitvoeringstijden voor RustProv en Vagrant bij het scenario met één
Linux-VM. De totale uitvoeringstijd wordt berekend door de gemiddelde starttijd
en de gemiddelde provisioningtijd bij elkaar op te tellen.

\section{Meting van meerdere Linux-VM's}

Voor het tweede scenario werden drie AlmaLinux-VM's gebruikt. De VM's werden één
voor één opgestart en vervolgens één voor één geprovisioned. In tegenstelling tot
het scenario met één Linux-VM wordt hier alleen de totale uitvoeringstijd van de
volledige Linuxopstelling gemeten.

De meting start wanneer het commando voor de volledige opstelling wordt
uitgevoerd en eindigt wanneer alle drie de VM's zijn opgestart en de bijhorende
provisioningsscripts succesvol na elkaar zijn uitgevoerd. De VM's werden tijdens
elke run in dezelfde volgorde verwerkt, zodat de uitvoeringsvolgorde tussen de
verschillende metingen gelijk bleef.

\begin{table}[htbp]
    \centering
    \caption{Totale uitvoeringstijd van meerdere Linux-VM's met RustProv en Vagrant.}
    \label{tab:meting-meerdere-linux}
    \begin{tabular}{lrrrr}
        \toprule
        \textbf{Tool} &
        \textbf{Run 1 (s)} &
        \textbf{Run 2 (s)} &
        \textbf{Run 3 (s)} &
        \textbf{Gemiddelde (s)} \\
        \midrule
        RustProv & -- & -- & -- & -- \\
        Vagrant  & -- & -- & -- & -- \\
        \bottomrule
    \end{tabular}
\end{table}

Tabel~\ref{tab:meting-meerdere-linux} toont de drie afzonderlijke metingen en
het gemiddelde voor de volledige Linuxopstelling. De gemeten tijd omvat zowel
het na elkaar opstarten van de drie VM's als het na elkaar uitvoeren van de
bijhorende provisioningsscripts.

\section{Meting van één Windows-VM}

Voor het derde scenario werd één Windows-VM gebruikt. De starttijd en de
provisioningtijd werden afzonderlijk gemeten voor RustProv en Vagrant. Voor elke
run werd dezelfde Windows-VM gebruikt en werd dezelfde beginsituatie voorzien.

De starttijd omvat het opstarten van de Windows-VM. De provisioningtijd omvat het
uitvoeren van het Windows-provisioningscript en het afronden van de bijhorende
configuratiestappen.

\begin{table}[htbp]
    \centering
    \caption{Meting van één Windows-VM met RustProv en Vagrant.}
    \label{tab:meting-een-windows}
    \begin{tabular}{llrrrr}
        \toprule
        \textbf{Tool} &
        \textbf{Onderdeel} &
        \textbf{Run 1 (s)} &
        \textbf{Run 2 (s)} &
        \textbf{Run 3 (s)} &
        \textbf{Gemiddelde (s)} \\
        \midrule
        RustProv & Start & -- & -- & -- & -- \\
        RustProv & Provisioning & -- & -- & -- & -- \\
        Vagrant  & Start & -- & -- & -- & -- \\
        Vagrant  & Provisioning & -- & -- & -- & -- \\
        \bottomrule
    \end{tabular}
\end{table}

Tabel~\ref{tab:meting-een-windows} toont de afzonderlijke metingen en de
gemiddelde uitvoeringstijden voor RustProv en Vagrant bij het scenario met één
Windows-VM. De gemiddelde waarden worden gebruikt om de uitvoeringstijd van
beide provisioners met elkaar te vergelijken.

\section{Meting van meerdere Windows-VM's}

Voor het vierde scenario werden drie Windows-VM's gebruikt: twee servers en één
client. De VM's werden tijdens elke run in dezelfde volgorde opgestart en
geprovisioned. Net zoals bij het meerledige Linuxscenario wordt hier alleen de
totale uitvoeringstijd van de volledige Windowsopstelling gemeten.

De meting start wanneer het commando voor de volledige Windowsopstelling wordt
uitgevoerd en eindigt wanneer alle drie de VM's zijn opgestart en de drie
Windowsprovisioningsscripts na elkaar succesvol zijn uitgevoerd.

\begin{table}[htbp]
    \centering
    \caption{Totale uitvoeringstijd van meerdere Windows-VM's met RustProv en Vagrant.}
    \label{tab:meting-meerdere-windows}
    \begin{tabular}{lrrrr}
        \toprule
        \textbf{Tool} &
        \textbf{Run 1 (s)} &
        \textbf{Run 2 (s)} &
        \textbf{Run 3 (s)} &
        \textbf{Gemiddelde (s)} \\
        \midrule
        RustProv & -- & -- & -- & -- \\
        Vagrant  & -- & -- & -- & -- \\
        \bottomrule
    \end{tabular}
\end{table}

Tabel~\ref{tab:meting-meerdere-windows} toont de drie afzonderlijke metingen en
het gemiddelde voor de volledige Windowsopstelling. De gemeten tijd omvat zowel
het na elkaar opstarten van de drie Windows-VM's als het na elkaar uitvoeren van
de drie Windowsprovisioningsscripts.

\section{Vergelijking tussen RustProv en Vagrant}
\label{sec:vergelijking-rustprov-vagrant}
